\documentclass[11pt,a4paper]{article}
\usepackage[spanish]{babel}
\usepackage[utf8]{inputenc}
\usepackage[T1]{fontenc}
\usepackage{amsmath,amssymb,booktabs,array,geometry,circuitikz,enumitem}
\geometry{margin=2.2cm}
\setlength{\parindent}{0pt}
\newcommand{\circuito}{%
\begin{center}
\begin{circuitikz}[american,scale=.88,transform shape]
\draw (0,0) to[battery1,l_=$V_1$] (0,3)
 to[R,l=$R_1$,i>^=$I_1$] (3,3) node[circ]{}
 to[R,l=$R_3$,i>^=$I_3$] (6,3) node[circ]{}
 -- (8,3) to[R,l=$R_5$] (8,0) -- (6,0) node[circ]{};
\draw (3,3) to[R,l=$R_2$,i>^=$I_2$] (3,0) node[circ]{} -- (0,0);
\draw (6,3) to[R,l=$R_4$] (6,0);
\draw (3,0) to[battery1,l=$V_2$] (6,0);
\end{circuitikz}
\end{center}}
\begin{document}
\begin{center}{\Large\bfseries Problema 2 --- Versión 1}\\[3pt]\end{center}
Considere el circuito compuesto de cinco resistencias y dos fuentes de voltaje como muestra la figura. Se tiene que:
\[
R_1=R_3=2\,\Omega,\quad R_2=4\,\Omega,\quad R_4=R_5=8\,\Omega,
\qquad V_1=16\,\mathrm V,\quad V_2=4\,\mathrm V.
\]. Entonces:
\circuito
\begin{enumerate}[label=\alph*)]
\item Determine la resistencia equivalente $R_{\mathrm{eq}}$ que se obtiene al reducir las resistencias $R_3$, $R_4$ y $R_5$. Indique las combinaciones en serie y en paralelo utilizadas (5 pts).
\item Aplicando las leyes de Kirchhoff, determine las corrientes $I_1$, $I_2$ e $I_3$, con los sentidos indicados en el esquema.(6 pts)
\item Determine la corriente, la tensión y la potencia disipada en $R_1$, $R_2$ y $R_{\mathrm{eq}}$. Complete la tabla(9 pts):
\begin{center}\begin{tabular}{|c|>{\centering\arraybackslash}p{2.8cm}|>{\centering\arraybackslash}p{2.8cm}|>{\centering\arraybackslash}p{2.8cm}|}\hline
Elemento & $I\,(\mathrm A)$ & $\Delta V\,(\mathrm V)$ & $P\,(\mathrm W)$\\\hline
$R_1$ &&&\\[12pt]\hline $R_2$ &&&\\[12pt]\hline $R_{\mathrm{eq}}$ &&&\\[12pt]\hline
\end{tabular}\end{center}
\end{enumerate}
\newpage
\begin{center}{\Large\bfseries Problema 2 --- Versión 1}\\[3pt]{\large Pauta }\end{center}
\textbf{a) Reducción de resistencias}\par
$R_4$ y $R_5$ están en paralelo: 
\[
\frac{1}{R_{45}}=\frac{1}{R_4}+\frac{1}{R_5}=\frac18+\frac18=\frac14
\quad\Rightarrow\quad \boxed{R_{45}=4\,\Omega}.\textcolor{red}{(2 pts)}
\]
En serie con $R_3$:
\[
\boxed{R_{\mathrm{eq}}=R_3+R_{45}=2+4=6\,\Omega}.\textcolor{red}{(3 pts)}
\] 
\textbf{b) Leyes de Kirchhoff}\par
Ley de nodos:
\begin{equation} I_1=I_2+I_3. \tag{1}\end{equation}
Malla 1:
\begin{equation} V_1-I_1R_1-I_2R_2=0
\quad\Rightarrow\quad 16-2I_1-4I_2=0. \tag{2}\end{equation}
Malla 2, con $R_{\mathrm{eq}}$:
\begin{equation} V_2+I_2R_2-I_3R_{\mathrm{eq}}=0
\quad\Rightarrow\quad 4+4I_2-6I_3=0. \tag{3}\end{equation}
La solución del sistema (1)--(3) es:
\[
\boxed{I_1=4\,\mathrm A\textcolor{red}{(2 pts)},\qquad I_2=2\,\mathrm A\textcolor{red}{(2 pts)},\qquad I_3=2\,\mathrm A\textcolor{red}{(2 pts)}.}
\]
Las tres corrientes resultan positivas en los sentidos señalados.

\medskip
\textbf{c) Corriente, voltaje y potencia disipada}\par
Usando $\Delta V=IR$ y $P=I^2R$:
\begin{center}
\begin{tabular}{lccc}
\toprule
Elemento & $I\,(\mathrm A)$ & $\Delta V\,(\mathrm V)$ & $P\,(\mathrm W)$\\
\midrule
$R_1=2\,\Omega$ & $4$ & $8$ & $32$\\
$R_2=4\,\Omega$ & $2$ & $8$ & $16$\\
$R_{\mathrm{eq}}=6\,\Omega$ & $2$ & $12$ & $24$\\
\bottomrule
\textcolor{red}{(1 punto cada una)}
\end{tabular}
\end{center}

\end{document}
